\documentclass[10pt,twocolumn]{article} 

\usepackage{simpleConference}
\usepackage{times}
\usepackage{apacite}
\usepackage{graphicx}
\usepackage{amssymb}
\usepackage{url}

\begin{document}

\title{A Sample Two-Column Paper Format}

\author{Carl Vogel \\
Computational Linguistics Group  \\
Trinity College Dublin \\
{\tt vogel@tcd.ie}}

\maketitle
\thispagestyle{empty}

\begin{abstract}
   This is a simple sample.
\end{abstract}

\section{Introduction}

This is the shell of a two-column paper.



\begin{figure}[hbt]
  \begin{center}
    \includegraphics[width=3.5in]{figure.eps}
  \end{center}
  \caption{\small TCD Logo.}
  \label{f:s}
\end{figure}


\begin{figure*}[hbt]
  \begin{center}
    \includegraphics[width=3.5in]{figure.eps}
  \end{center}
  \caption{\small This is a two column version of the TCD logo.}
  \label{f:t}
\end{figure*}

\section{Readings}

We read a work of \citeA{Kuh70}, and also of
\citeauthor{Quine2d53}.  Some texts were foundational and formal
\cite{Gar79}.  Some texts were foundational and philosophical
\cite{Descartes1637}.  We considered \citeA{Morton76} as a marvel of
conciseness.\footnote{Notice how the sentences in this section
are grammatical:  some include the author's name as a word in the
sentence and some make the textual citation wholly parenthetical.}

\section{Conclusion}

This is more terse than a proper paper would be.\footnote{Note that
three bibliographic files require naming: a bibliography, a bibligraphy
style, and a bibliographic package (apacite appears in two of those; 
the author defines the third).}

\bibliographystyle{apacite}
\bibliography{bibliography}
\end{document}
